\section{总结与延伸}

\subsection{本期核心结论}

\begin{enumerate}
\item Agent 不是新概念，AI 从早期就有 Agent 想象；新变化是 LLM 让语言成为通用工具接口。
\item Language Agent 继承了 semantic parsing 的可执行性，但通过 LLM 获得更强泛化。
\item OpenClaw Moment 类似 ChatGPT Moment，都是已有能力被产品形态放大后触发社会共识。
\item Browser、desktop、mobile、GUI、CLI、API、coding 的边界正在消融，coding 是 digital world 的基础表达层。
\item Agent 的核心瓶颈可统一为 continual learning 与 world model：从经验中学习，形成可靠专家能力。
\item 2026 年 Agent 竞争不只是模型竞争，也是产品、部署、成本、可靠性和组织执行竞争。
\end{enumerate}

\subsection{开放问题}

\begin{itemize}
\item Continual learning 的主流路线会是基于 world model 的学习，还是更工程化的记忆/检索/工作流系统？
\item Universal digital agent 会先在 coding/productivity 中成熟，还是在消费端应用中爆发？
\item GUI、CLI、API、coding 的边界会在模型层统一，还是长期保留多套产品入口？
\item Forward-deployed engineer 模式会是过渡形态，还是 Agent 企业落地的长期常态？
\end{itemize}

\subsection{拓展阅读}

\begin{itemize}
\item \textit{Artificial Intelligence: A Modern Approach}：理解早期 AI 与 Agent 定义的经典入口。
\item ReAct、Toolformer、AutoGPT、Mind2Web、LM Planner 等工作：理解 Language Agent 最近三年的关键节点。
\item Computer Use Agent、Web Agent、GUI Agent、Coding Agent 相关 benchmark：理解 universal digital agent 的评估难点。
\end{itemize}

\begin{importantbox}{最后的压缩}
Agent 的终点不是“更会聊天”，而是“更会在环境中积累经验并完成目标”。OpenClaw Moment 的意义，是让行业看见这种可能性；接下来的难点，是把可能性变成可靠、低成本、可部署的系统。
\end{importantbox}

\end{document}
